\chapter{自动G点：什么时候真有一张固定网络}\label{ch:networks}
\section{所有有理输入可解，竟然能推出同一个词}
若固定一个平均词，对所有输入都达到全局均值，它的矩阵必为 $J_n/n$。条目属于 $\Z[1/k]$，因此 $n$ 的素因子都必须整除 $k$。这是固定网络的必要条件，24平均27元说明它还不充分。

更有用的是反向的一条有限维论证。若每个有理零和向量都可被某个词消去，取零和空间的一组有理基。先消去第一基向量；再为第二基向量的\emph{当前像}选一个消去词，接在后面；已经为零的像仍为零。有限次后，整个零和空间都被一个词消掉。平均词固定常值方向，所以这个词恰为 $J_n/n$。\src{S01}\src{S05}

因此
\[
\text{固定全输入网络存在}
\iff\bigl[\rad(n)\mid k\text{ 且 }\mathcal P(k,n)\bigr].
\]
右边并没有省掉 $\mathcal P(k,n)$；自动G合法仍须证明可达。

\section{六平均九元：四步里每一步在干什么}
先将任意六项平均为 $a$，另三项均值记为 $c$。第二步取三份 $a$ 与那三项，得到六份 $b=(a+c)/2$，另外三份仍为 $a$。
全局均值是 $\mu=(2a+c)/3$。第三步取两份 $a$、四份 $b$，均值恰为 $\mu$；剩下三项为 $a,b,b$，和为 $3\mu$。第四步取它们与三份 $\mu$，九项全等。\src{S09}

例如第一组均值 $a=3$、外部均值 $c=9$，第二步得 $b=6$，第三步得 $(2\cdot3+4\cdot6)/6=5$，恰为九元均值。剩余 $3,6,6$ 加三份五，也平均成五。整个词不依赖这些具体数值。

\begin{theorem}[四步网络族]\label{thm:four-network}
设 $k=ur$，$u\ge1$，且 $u+1\mid r$。在 $n=k+r$ 个位置上，存在固定四步 $k$ 平均网络。
\end{theorem}
\begin{proof}
第一步平均 $k$ 项为 $a$，余下 $r$ 项均值为 $c$。第二步取 $k-r$ 份 $a$ 与全部外部位置，得 $k$ 份 $b=((u-1)a+c)/u$，留下 $r$ 份 $a$。全局均值为 $\mu=(a+ub)/(u+1)$。
令 $h=ur/(u+1)=kr/(k+r)$，整除条件使它为整数。第三步取 $h$ 份 $a$、$k-h$ 份 $b$，均值为 $\mu$。剩余 $r$ 项总和必为 $r\mu$，第四步将它们与 $k-r$ 份 $\mu$ 平均即可。
\end{proof}
取 $k=m(m-1),r=m$，得到 $n=m^2$ 的网络：$2\to4,6\to9,12\to16,20\to25,30\to36,42\to49$ 等。若 $m$ 或 $m-1$ 为素数，它还与上一讲的下界相接，给出 $M(m(m-1))=m^2$。两者都合数时只能给一个上界，例如 $77\le M(72)\le81$。\src{S01}

固定全输入网络少于四步不可能：两步后，若新行纳入第一组外的位置，该输入位置的系数为 $1/k\ne1/n$；若没纳入，则支撑不足。未动行也不是全均值行。因此两步后没有任何一行正确，第三步又必留下 $n-k>0$ 行不动，不能全体正确。这是固定词下界，不是每个具体输入的最少步数。

\section{十平均十六元：用收缩误差配上最后的单点}
十六元不在上述四步族里，但有一张六步网络。更一般地，取
\[n=4^b,\qquad h=(n-1)/3,\qquad k=2h,\qquad b\ge1.
\]
先平均 $2h$ 个位置为 $a$，再平均 $h-1$ 份 $a$ 与全部其余 $h+1$ 个位置，得 $2h$ 份 $b_0$ 和 $h+1$ 份 $a$。留下一个 $a$ 作单点，其余整理为三个等大 $h$ 块，块值是 $(a,b_0,b_0)$。\src{S12}

现在每次合并两个 $h$ 块做一次 $k$ 平均，等于在三块值上做二平均。依次选块对 $(0,1),(1,2),(2,0)$ 循环。令 $a$ 的系数为一、$b_0$ 的系数为零，第 $j$ 次产生的新系数是
\[c_j=\frac{1-(-1/2)^j}{3},\qquad c_0=0,c_1=1/2.
\]
递推 $c_j=(c_{j-1}+c_{j-2})/2$ 可直接验证该公式。写 $m=2b$，在 $j=m-1$ 时
\[c_{m-1}=\frac{1+2^{1-m}}3=\frac{h+1}{3h+1},
\]
正好是全局均值里 $a$ 的系数。于是已有 $2h$ 份全局均值，最后将未达标的 $h+1$ 项与 $h-1$ 个正确位置平均，全部完成。总步数 $2+(m-1)+1=2b+2$。

对十平均十六元，后三次块混合依次得到 $(a+b_0)/2,(a+3b_0)/4,(3a+5b_0)/8$；最后一项就是六份 $a$ 与十份 $b_0$ 的总体均值。故 $M(10)=16$；结合十七元控制器，$H(10)=17$。六步没有被证明最短。

\section{张量网络：按坐标方向轮流平均}
若 $k=\prod p^{a_p}$，位置数 $B=\prod p^{b_p}$ 满足每个 $b_p\ge a_p$，把位置标为各素数进位坐标的笛卡尔积。每层选每个 $p$ 的 $a_p$ 条坐标轴，在对应的 $k$ 元纤维上平均。轮换轴，覆盖全部坐标，得到全体平均。\src{S05}

所需层数可取 $L=\max_p\lceil b_p/a_p\rceil$，操作数为 $LB/k$。证明只用不同坐标方向的全平均算子彼此交换。注意这是充分构造，不是所有固定网络的分类；六平均九元中 $6\nmid9$，仍有刚才的四步网络。

在 $3^b$ 个位置上，把九元六平均网络放到两条三进坐标轴组成的每个纤维，轮换覆盖所有轴，得到 $b\ge2$ 时的固定网络，长度
\[4\cdot3^{b-2}\lceil b/2\rceil.
\]
这些网络会作为下一讲块均值模拟的实际子程序。
